\begin{tikzpicture}[x=1cm,y=1cm,every node/.append style={anchor=center,inner sep=.3333em}]
\draw[GeoA,line width=1.2pt] (-.4,-.2)--(6.4,3.2);
\coordinate (P) at (2.2,2.5);
\coordinate (Q) at (3.25,1.625);
\draw[GeoB,line width=.8pt,fill=GeoB!10] (P) ellipse(1.15 and .32);
\fill[Ink] (P) circle(2.2pt) node[pLabel,anchor=south,yshift=2.4mm]{$(X_i,Y_i)$};
\draw[-{Latex[length=2mm]},Classical,line width=.9pt] (P)--(Q);
\fill[Classical] (Q) circle(2pt) node[pNote,anchor=north west,text=Classical]{punto corretto};
\draw[Muted,dashed] (P)--(2.2,1.1);
\node[pNote,text=Muted,anchor=east] at(2.15,1.75){residuo verticale};
\node[pNote,text width=7.4cm] at(3,-1.2){L’ellisse rappresenta gli errori del campione: ampia in orizzontale, stretta in verticale. York sposta il punto sulla retta nella direzione più compatibile con i suoi errori e pesa ciascun residuo con $W_i$.};
\end{tikzpicture}
